\documentclass[11pt]{article}
\usepackage[margin=1.25in]{geometry}
\usepackage[T1]{fontenc}
\usepackage{hyperref}

\title{ProvBench-RAG: A Benchmark Resource for Provenance-Sensitive Retrieval Under Near-Duplicate Source Confusion}
\author{[Authors --- Resource track, single-blind]}
\date{}

\begin{document}
\maketitle
\thispagestyle{empty}

\noindent\textbf{Track:} CIKM 2026 Resource \\
\noindent\textbf{Keywords:} benchmark, retrieval, provenance, RAG, dataset, evaluation

\vspace{0.75em}
\noindent\rule{\textwidth}{0.4pt}
\vspace{0.75em}

\begin{abstract}
\noindent
Retrieval-augmented systems are often evaluated on whether some passage supports an answer, but deployed pipelines must also choose the correct document instance among near-duplicates (mirrors, stale snapshots, scraped copies, jurisdiction variants).
We release ProvBench-RAG, an open benchmark resource comprising JSONL schemas, three reproducible tracks (1,500 synthetic clusters, 140 Wikipedia revision clusters, 1,640 pooled queries), deterministic evaluation code, and reference baseline rankings.
The resource exposes preferred-source accuracy (PSA@1) and composite ProvScore alongside recall and nDCG.
On the synthetic track, lexical BM25 achieves PSA@1 of 0 under planted decoys while a bundled provenance reranker reaches 0.75; Wikipedia and pooled tracks quantify transfer and heterogeneous-index stress.
All data, metrics, and paper tables regenerate from public scripts and frozen annotations in the repository.
\end{abstract}

\end{document}
